\section{结果}

\subsection{原序列覆盖与 SWE 同降，SWE 趋势对同化断点敏感}

多年平均覆盖率呈强纬向梯度，从 $0$--$30^\circ$N 的 $0.78\%$ 增至 $65$--$90^\circ$N 的 $97.58\%$（图~\ref{fig:climatology}，表~\ref{tab:spatial}）。北部高覆盖区内，SWE 的空间差异仍很大。全分析域 43 个冬季平均覆盖面积为 $39.33\times10^6~\mathrm{km^2}$，平均 SWE 为 68.17~mm，平均 T2m 为 $0.55\,^\circ$C。

\fig{fig02_climatology}{fig:climatology}{\textbf{冬季积雪覆盖、雪水当量和气温的多年平均。} \textbf{a}，SCF 比例。 \textbf{b}，SWE（mm）。 \textbf{c}，T2m（$^\circ$C）。每格点对 1980--2022 年 43 个完整 DJF 冬季等权平均。固定分析域内有 11\,844 个格点。色条延伸端表示超出显示范围，空白表示分析域外。}

1980--2022 年间，SCA 每十年减少 $0.575\times10^6~\mathrm{km^2}$（\CI\ $-0.807$ 至 $-0.344$，$p=1.05\times10^{-5}$），相当于 SCR 每十年减少 0.586 pp。同期平均 SWE 每十年减少 2.233~mm，T2m 每十年增加 $0.318\,^\circ$C（图~\ref{fig:global}，表~\ref{tab:global}，各 $n=43$）。MK 和 Sen 斜率支持同一方向。储量情景 A 和 B 的斜率分别为 $-219.5$ 和 $-208.1~\mathrm{km^3}$ 每十年。剔除任何冬季 SWE 曾超过 1\,000~mm 的 224 个格点后，情景 A 的均值由 6\,700 降为 2\,808~km$^3$，斜率仍为 $-231.5~\mathrm{km^3}$ 每十年（\CI\ $-287.2$ 至 $-175.7$）。高值格点主要影响储量基线。

\begin{table}[!htbp]\centering\small
\caption{\textbf{全分析域冬季指标的线性趋势。} 1980--2022 年，$n=43$ 个 DJF 冬季。各斜率以每十年计。区间和双侧 $p$ 值采用 HAC(3)，$n_{\rm eff}$ 为基于一阶残差自相关的描述性近似。}\label{tab:global}
\begin{tabular}{lrrrr}
\toprule
指标 & 斜率 [95\% CI] & $p_{\rm HAC}$ & $R^2$ & $n_{\mathrm{eff}}$ \\
\midrule
SCA ($10^6$ km$^2$) & $-0.575\;[-0.807,\ -0.344]$ & $1.05\times 10^{-5}$ & $0.387$ & $43.0$ \\
SCR (pp) & $-0.586\;[-0.821,\ -0.350]$ & $1.05\times 10^{-5}$ & $0.387$ & $43.0$ \\
平均 SWE (mm) & $-2.233\;[-2.873,\ -1.594]$ & $1.39\times 10^{-8}$ & $0.652$ & $13.4$ \\
T2m ($^\circ$C) & $0.318\;[0.193,\ 0.443]$ & $7.27\times 10^{-6}$ & $0.460$ & $40.4$ \\
$V_A$ (km$^3$) & $-219.5\;[-282.3,\ -156.6]$ & $1.39\times 10^{-8}$ & $0.652$ & $13.4$ \\
$V_B$ (km$^3$) & $-208.1\;[-269.1,\ -147.1]$ & $2.36\times 10^{-8}$ & $0.641$ & $13.7$ \\
\bottomrule
\end{tabular}

\end{table}

\fig{fig03_global_series}{fig:global}{\textbf{全分析域积雪覆盖和雪水指标共同下降。} \textbf{a}，积雪覆盖面积。 \textbf{b}，面积加权平均 SWE。 \textbf{c}，储量情景 A 与 B，分别假定 SWE 为全格点平均或雪覆盖区条件平均。 \textbf{d}，面积加权 T2m。标记为 43 个 DJF 冬季原值，线为 OLS 拟合，带为 HAC(3) 拟合均值的逐点 95\% 置信区间。}

加入 2004 年 1 月的水平项后，平均 SWE 斜率缩小至 $-0.760$~mm 每十年（\CI\ $-2.045$ 至 $0.525$，$p=0.239$），两种储量斜率的区间也均跨零。SCR 斜率仍为 $-0.686$ pp 每十年（\CI\ $-1.102$ 至 $-0.269$，$p=0.0019$），T2m 仍为 $+0.385\,^\circ$C 每十年（\CI\ $0.186$ 至 $0.585$，$p=0.00036$，均 $n=43$）。同化变更前后的分段拟合给出方向不同的 SWE 斜率。这个敏感性检验降低了将原 SWE 斜率解释为连续气候趋势的可信度。

SCA 与平均 SWE 的原序列相关为 $r=0.484$，去趋势后为 $r=-0.038$。两者标准化斜率差为 $+0.148$ 个标准差每十年（5 冬季区块 \CI\ $-0.167$ 至 $0.466$，$p=0.367$，$n=43$）。用 $V_B$ 替代平均 SWE 后结论一致，剔除高 SWE 格点后去趋势相关仍接近零。共同下降可解释原序列的部分相关。年际异常缺乏线性同步的证据，标准化长期斜率差仍未得到明确支持。

\subsection{覆盖率下降在中纬度集中，SWE 下降在较高纬度更快}

格点 OLS 斜率的正负分布在三变量间不同。负 SCF 斜率占分析域面积的 41.15\%，负 SWE 斜率占 58.88\%，正 T2m 斜率占 91.39\%。按 HAC 后逐变量做 BH 校正，SCF、SWE 和 T2m 分别有 1.71\%、26.01\% 和 40.14\% 的格点通过 $q<0.05$，对应面积 1.81\%、21.52\% 和 46.42\%（图~\ref{fig:trends}）。常规 OLS 或 MK 后再做同样校正时，无 SCF 格点通过。覆盖率下降在分析域聚合后较清晰，具体显著格点依赖检验方法。相邻斜率明显相似，SCF、SWE 和 T2m 的 Moran $I$ 分别为 0.792、0.394 和 0.817（相对随机交换斜率位置，均 $p=0.001$）。该检验描述估计斜率的空间聚集。SCF 和 SWE 点估计反号的格点占面积的 10.68\%，没有格点在两者均通过主校正时呈相反方向。

\fig{fig04_grid_trends}{fig:trends}{\textbf{三变量趋势的空间分布。} \textbf{a}，SCF（pp 每十年）。 \textbf{b}，SWE（mm 每十年）。 \textbf{c}，T2m（$^\circ$C 每十年）。斜率按 1980--2022 年 43 个 DJF 冬季逐格 OLS 拟合。点标出逐变量经 BH 校正后 HAC(3) $q<0.05$ 的格点。色标在对应变量的综合图中沿用，端部延伸表示超出显示范围。}

欧洲窗口 SCF 下降 $1.560$ pp 每十年，北美洲窗口估计下降 $0.369$ pp 每十年。两窗口的直接斜率差为 $-1.192$ pp 每十年（5 冬季区块 \CI\ $-2.029$ 至 $-0.364$，$q=0.016$，各 $n=43$）。三个窗口的 SWE 均下降，但其窗口间斜率差均未通过组内校正。欧洲窗口的 T2m 增幅大于亚洲窗口，差为 $0.303\,^\circ$C 每十年（\CI\ $0.148$ 至 $0.458$，$q=0.00075$）。调整经纬度边界后，三变量斜率点估计的方向和排序一致。SCA 与 SWE 标准化斜率差的窗口间对比未通过校正，现有数据支持区域覆盖率趋势有差异，对覆盖--储量脱耦程度的区域排序证据较弱（图~\ref{fig:windows}）。

\fig{fig05_approximate_geographic_windows}{fig:windows}{\textbf{代表窗口的积雪、气温和覆盖--储量同步性。} \textbf{a--c}，1980--2022 年 43 个 DJF 冬季的面积加权 SCF、SWE 与 T2m，虚线为 OLS 拟合。 \textbf{d--f}，每十年斜率及 HAC(3) 95\% CI。 \textbf{g}，SCA 与平均 SWE 的标准化异常差 $D_t$。 \textbf{h}，以平均 SWE、$V_A$、$V_B$ 分别计算的 $D_t$ 斜率及 5 冬季区块自助区间。 \textbf{i}，去趋势前后 SCA 与 SWE 的 Pearson 相关。各区域均为近似经纬度窗口。}

$30$--$50^\circ$N 的 SCF 下降最快，每十年减少 $1.347$ pp（\CI\ $-1.952$ 至 $-0.743$，$q=2.18\times10^{-4}$）。$50$--$65^\circ$N 和 $65$--$90^\circ$N 的 SWE 分别每十年减少 $6.066$ 和 $6.343$~mm，两个区间均低于零（表~\ref{tab:spatial}，图~\ref{fig:latitude}）。最北带平均 SCF 达 $97.58\%$，覆盖变化仅 $-0.019$ pp 每十年，其 SWE 仍有明显变化。低纬带有效面积达 $35.70\times10^6~\mathrm{km^2}$，但平均 SCF 仅 $0.78\%$，平均 SWE 仅 $0.90$~mm。在此低雪基线上，SCA 与平均 SWE 标准化趋势差为 $-0.639$ 个标准差每十年（区块 \CI\ $-0.985$ 至 $-0.270$，$q=0.005$）。该差异反映各自波动尺度上的变化，不能直接外推为大范围储量变化。

\begin{table}[!htbp]\centering\small
\caption{\textbf{代表窗口与纬度带的有效面积和冬季趋势。} 每组 $n=43$，时期为 1980--2022 年 DJF。斜率单位均为每十年，方括号内为 HAC(3) 95\% CI。窗口内面积按固定有效格点计算。组内校正结果和斜率对比见源数据。}\label{tab:spatial}
\setlength{\tabcolsep}{4pt}
\begin{tabular}{lrrrr}
\toprule
分组 & 面积 ($10^6$ km$^2$) & SCF (pp) & SWE (mm) & T2m ($^\circ$C) \\
\midrule
欧洲窗口 & 12.90 & \shortstack[r]{ $ -1.560 $\\$[-2.393,\ -0.728]$ } & \shortstack[r]{ $ -3.877 $\\$[-5.194,\ -2.560]$ } & \shortstack[r]{ $ 0.496 $\\$[0.175,\ 0.817]$ } \\
亚洲窗口 & 35.15 & \shortstack[r]{ $ -0.802 $\\$[-1.058,\ -0.546]$ } & \shortstack[r]{ $ -2.741 $\\$[-3.971,\ -1.511]$ } & \shortstack[r]{ $ 0.193 $\\$[-0.030,\ 0.415]$ } \\
北美洲窗口 & 21.76 & \shortstack[r]{ $ -0.369 $\\$[-0.911,\ 0.173]$ } & \shortstack[r]{ $ -3.050 $\\$[-4.330,\ -1.771]$ } & \shortstack[r]{ $ 0.348 $\\$[0.171,\ 0.526]$ } \\
$0$--$30^\circ$N & 35.70 & \shortstack[r]{ $ -0.061 $\\$[-0.098,\ -0.024]$ } & \shortstack[r]{ $ 0.147 $\\$[-0.042,\ 0.335]$ } & \shortstack[r]{ $ 0.306 $\\$[0.246,\ 0.367]$ } \\
$30$--$50^\circ$N & 32.04 & \shortstack[r]{ $ -1.347 $\\$[-1.952,\ -0.743]$ } & \shortstack[r]{ $ -1.153 $\\$[-1.683,\ -0.623]$ } & \shortstack[r]{ $ 0.315 $\\$[0.160,\ 0.470]$ } \\
$50$--$65^\circ$N & 21.54 & \shortstack[r]{ $ -0.558 $\\$[-0.853,\ -0.262]$ } & \shortstack[r]{ $ -6.066 $\\$[-7.859,\ -4.273]$ } & \shortstack[r]{ $ 0.267 $\\$[-0.007,\ 0.542]$ } \\
$65$--$90^\circ$N & 9.00 & \shortstack[r]{ $ -0.019 $\\$[-0.031,\ -0.007]$ } & \shortstack[r]{ $ -6.343 $\\$[-8.746,\ -3.941]$ } & \shortstack[r]{ $ 0.495 $\\$[0.180,\ 0.809]$ } \\
\bottomrule
\end{tabular}

\end{table}

\fig{fig06_latitude_response}{fig:latitude}{\textbf{覆盖率与 SWE 的长期下降在不同纬度带最强。} \textbf{a--c}，四个纬度带的 OLS 每十年斜率和 HAC(3) 95\% CI。 \textbf{d}，控制年份的温度--积雪 Pearson 偏相关。 \textbf{e,f}，去趋势 SCF 与 SWE 对去趋势 T2m 的回归斜率及 HAC(3) 95\% CI，实心点表示组内校正后 $q\leq0.05$。每组含 1980--2022 年 43 个 DJF 冬季。}

\subsection{阈值雪线的分析域中位位置北移}

在 1980--1989 和 2013--2022 年平均场均可识别 $\theta=0.50$ 交叉的 174 条经度上，晚期相对早期的纬度位移中位数为 $+0.366^\circ$（10 经度空间区块 \CI\ $0.207^\circ$ 至 $0.593^\circ$，四分位距 $0.105^\circ$ 至 $1.341^\circ$）。该统计先配对再汇总。三个代表窗口的对应位移分别为 $+0.450^\circ$、$+0.419^\circ$ 和 $+0.263^\circ$（图~\ref{fig:boundary}）。在阈值 0.25 和 0.75 下，全分析域中位位移均向北。窗口排序随阈值改变，等值线叠加图也显示不同经度可有相反方向（附录表~\ref{tab:boundary}）。

\fig{fig07_threshold_snowline}{fig:boundary}{\textbf{早晚期覆盖场与 SCF 阈值雪线。} \textbf{a,b}，1980--1989 和 2013--2022 年各 10 个 DJF 冬季平均 SCF。 \textbf{c}，晚减早 SCF 差（pp），虚线和实线分别为两期 $\theta=0.50$ 等值线。 \textbf{d}，阈值 0.25、0.50 和 0.75 下全域和各代表窗口的配对经度位移中位数，区间为 10 经度区块 95\% CI。边界表示二维覆盖位置。}

以逐冬季经度中位数拟合，$\theta=0.50$ 的全分析域边界每十年向北移动 $0.393^\circ$（\CI\ $0.111^\circ$ 至 $0.675^\circ$，$p=0.0075$，$q=0.0129$，$n=43$）。逐冬季可用经度数为 173--193。固定 160 条全时段均可用的经度后，斜率为 $+0.403^\circ$ 每十年（\CI\ $0.111^\circ$ 至 $0.694^\circ$）。加入 2004 年水平项后，以可用经度拟合的斜率仍为正（$+0.698^\circ$ 每十年，\CI\ $0.166^\circ$ 至 $1.231^\circ$）。另外两个 SCF 阈值的年度斜率也为正。改变孤立斑块规则时，主阈值斜率为每十年 $+0.413^\circ$ 至 $+0.551^\circ$。欧洲窗口年度趋势为 $+0.487^\circ$ 每十年（$q=0.0032$），亚洲窗口的证据在固定经度或多重校正下减弱，北美洲窗口未检出明确的年度线性趋势（$q=0.852$）。后者的早晚期配对位移不能单独证明持续线性移动（图~\ref{fig:boundaryannual}）。

\fig{fig08_snowline_series}{fig:boundaryannual}{\textbf{主阈值边界的年度中位纬度和有效经度。} 左列为 SCF $=0.50$ 年度边界纬度中位数、OLS 线及 HAC(3) 逐点 95\% 拟合区间。右列为同冬季的有效经度数。各行依次为全分析域、欧洲、亚洲及北美洲代表窗口。每个趋势使用 1980--2022 年 43 个 DJF 冬季。}

剔除过渡冬季并加入水平项后，主阈值全域斜率仍向北（$+0.696^\circ$ 每十年，$n=42$），对 12 个阈值与区域组合校正后证据减弱（$q=0.090$）。

\subsection{去趋势后的温度关联因指标和纬度而异}

全分析域去趋势 SCF 与 T2m 的偏相关为 $-0.679$，斜率为 $-1.455$ pp/$^\circ$C（\CI\ $-1.922$ 至 $-0.989$，HAC $p=1.76\times10^{-7}$，$R^2=0.461$，$n=43$）。SWE 与 T2m 的原始相关为 $-0.439$，去趋势后为 $+0.251$，残差斜率为 $+1.191$ mm/$^\circ$C（\CI\ $0.277$ 至 $2.106$，HAC $p=0.012$，$R^2=0.063$）。SWE 残差仍有较强一阶相关（约 $0.495$）。逐冬季删除所得 SWE 斜率保持为正，范围 $1.010$--$1.698$ mm/$^\circ$C。常规 OLS 的 $p=0.109$。加入 2004 年水平项后，SWE 斜率为 $+0.928$ mm/$^\circ$C（\CI\ $-0.017$ 至 $1.873$，$p=0.054$），SCF 斜率仍为 $-1.455$ pp/$^\circ$C（\CI\ $-1.942$ 至 $-0.967$）。SWE 当季残差正关联的解释力较低，推断对模型设定敏感（图~\ref{fig:temperature}）。

SCF 的温度斜率在 $30$--$50^\circ$N 最负，为 $-2.254$ pp/$^\circ$C（\CI\ $-3.066$ 至 $-1.441$，$q=6.73\times10^{-6}$）。$65$--$90^\circ$N 的斜率为 $+0.023$ pp/$^\circ$C（\CI\ $0.009$ 至 $0.038$，$q=0.0027$），幅度很小。纬度带间 SCF 斜率差得到联合检验支持（5 冬季区块 $q=0.0005$）。各纬度带的 SWE 斜率均未通过组内校正，SWE 的联合检验随区块长度变化（3、5、7 冬季区块的 $q$ 分别为 0.112、0.0245、0.0233）。亚洲和北美洲窗口的 SWE 斜率则分别为 $+2.106$ 和 $-1.762$ mm/$^\circ$C（$q=0.0031$ 与 $0.0091$），窗口间联合检验支持差异（5 冬季区块 $q=0.0005$）。这些系数对应区域聚合的统计关联。

在 80 项跨冬季滞后检验中，校正后无正滞后关系达到 $q<0.05$。在 $k=-2$ 的北美洲窗口和低纬带出现两个 SCF 信号，方向均为积雪领先气温（$q=0.0225$、0.0221）。这些信号出自探索性比较，现有序列对前冬气温预测后冬积雪的支持有限。

\fig{fig09_temperature_lags}{fig:temperature}{\textbf{当季温度关联与跨冬季滞后。} \textbf{a,d}，全分析域的去趋势异常与 OLS 拟合。 \textbf{b,e}，全域、近似窗口与纬度带控制年份的偏相关（$n=43$），区间为 5 冬季配对区块 95\% CI，实心点表示该组回归 HAC 检验经 BH 校正后 $q\leq0.05$。 \textbf{c,f}，全分析域跨冬季相关（$k=-2$ 至 $2$，$n=41$--43），误差线为 5 冬季区块区间，填充按全 80 项滞后检验校正。正滞后表示温度领先。}

\fig{fig10_synthesis}{fig:synthesis}{\textbf{冬季覆盖、SWE、气温趋势与地理积雪边界变化。} \textbf{a--c}，1980--2022 年 43 个 DJF 冬季的格点 SCF、SWE 和 T2m 斜率，单位分别为 pp、mm 和 $^\circ$C 每十年。点为逐变量 BH 校正后 HAC(3) $q<0.05$。 \textbf{d}，2013--2022 年相对 1980--1989 年的 SCF 差值（pp）及两期 SCF $=0.50$ 等值线。各期包含 10 个 DJF 冬季，空白表示分析域外。}

阈值 0.25 和 0.75 的早晚边界叠加见附录图~\ref{fig:thresholdsens}。逐经度边界的样本数、位移分位数与敏感性结果保存在 \nolinkurl{output/csv/snowline_paired_shifts.csv}。
